\section{From pivots to a gluing lattice}
\label{glue:section}

Fix the operator $a$ of Proposition~\ref{op:operator}. We first pass
from its finite-block pivots to actual kernel and cokernel vectors in
$\Phi$. The finite-crossing property makes the resulting quotient
Hausdorff and controls its compact open subspaces. We then pair the
two sorts of defects in a discrete lattice and glue it to the cone of
$a$. The resulting object will satisfy the first two conditions of
the transfer criterion; exclusion from $\Thick(T')$ will be the next
section's task.

For $J\in\Z$, put
\begin{equation}
 F_J\Phi=w^Jk[[w^{-1}]],\qquad
 \mathcal Q=\Phi/a\Phi,\qquad
 \mathcal K=\ker(a:\Phi\longrightarrow\Phi).
 \label{glue:spaces}
\end{equation}
The $F_J\Phi$ are compact open subspaces for the Laurent-series
topology, with $\Phi=\bigcup_JF_J\Phi$ and
$\bigcap_JF_J\Phi=0$. Give $\mathcal K$ the subspace topology and
$\mathcal Q$ the quotient topology. The \emph{leading position} of
a nonzero series is its greatest exponent of $w$.

The operator $a$ is continuous and preserves every $F_J\Phi$.
Consequently, for $L\le J$, its matrix block on indices
$L,\ldots,J$ is the induced map on the finite-dimensional space
$F_J\Phi/F_{L-1}\Phi$. This remains true for inputs with infinitely
many nonzero coefficients below $L$: their tails contribute nothing
to the displayed rows. The rank of this block is the
$\mathsf R(L,J)$ used in Lemma~\ref{op:matching}.

\begin{lemma}[Pivots and the pole topology]
\label{glue:pivot-topology}
An integer $j$ is an unmatched column of $a$ if and only if there is
a vector in $\mathcal K$ with leading term $w^j$. An integer $l$ is
a matched row if and only if it is the leading position of a nonzero
vector in $a\Phi$. Moreover, for every $J\in\Z$,
\begin{equation}
 \dim_k\frac{a\Phi\cap F_J\Phi}{a(F_J\Phi)}
   =\#\{(l,j)\text{ matched}:l\le J<j\}<\infty.
 \label{glue:crossing-quotient}
\end{equation}
The image $a\Phi$ is closed. Thus $\mathcal Q$ and $\mathcal K$ are
Hausdorff locally compact $k$-vector spaces, with compact open
subspaces given respectively by the images of $F_J\Phi$ and by
$\mathcal K\cap F_J\Phi$.
\end{lemma}

\begin{proof}
Suppose column $j$ is unmatched. For every lower cutoff $L\le j$,
its rank increment in the block $[L,j]$ is zero. Hence
\[
 S_L=\{f\in w^j+F_{j-1}\Phi:a(f)\in F_{L-1}\Phi\}
\]
is nonempty: a solution of the finite coefficient equations extends
by a zero tail. Continuity of $a$ makes $S_L$ closed in the compact
space $w^j+F_{j-1}\Phi$. Lowering $L$ adds equations, so these
sets are nested. Compactness gives a point in their intersection,
and its image lies in $\bigcap_LF_{L-1}\Phi=0$. This is an actual
kernel vector with leading term $w^j$. Conversely, such a vector
makes column $j$ dependent in every finite block, so the column is
unmatched. In particular,
\begin{equation}
 \dim_k\frac{a(F_j\Phi)}{a(F_{j-1}\Phi)}
 =
 \begin{cases}
 1,&\text{if column $j$ is matched},\\
 0,&\text{if column $j$ is unmatched}.
 \end{cases}
 \label{rew:glue:image-increment}
\end{equation}
Indeed, $F_j\Phi/F_{j-1}\Phi$ is generated by $w^j$, and
$a(w^j)\in a(F_{j-1}\Phi)$ is equivalent to the existence of the
kernel vector just constructed.

If row $l$ is matched, then at some finite column cutoff its
addition increases rank. A finite column combination therefore
has nonzero coefficient in that row and zero coefficients above
it. The combination is an actual Laurent polynomial input, and
triangularity makes all rows above the finite block vanish too.
Conversely, suppose $a(f)$ has leading position $l$. Write
$f=f_{\ge l}+f_{<l}$, where $f_{\ge l}$ is a Laurent polynomial
and $f_{<l}\in F_{l-1}\Phi$. Since $a$ preserves $F_{l-1}\Phi$,
the infinite tail contributes nothing on rows at least $l$.
The finite polynomial $f_{\ge l}$ therefore exhibits the same
row increment in a finite block, so row $l$ is matched.

For $J'>J$, equation~\eqref{rew:glue:image-increment} shows that
the dimension of $a(F_{J'}\Phi)/a(F_J\Phi)$ is the number of
matched columns in $(J,J']$. The map from this quotient to
$F_{J'}\Phi/F_J\Phi$ induced by projection has rank
$\mathsf R(J+1,J')$: the images of input terms at or below $J$
lie in $F_J\Phi$ and vanish under the projection. The kernel of
this induced map is
\[
 \frac{a(F_{J'}\Phi)\cap F_J\Phi}{a(F_J\Phi)}.
\]
Subtracting the projected rank therefore counts exactly the matches
with $l\le J<j\le J'$. Every Laurent input belongs to some
$F_{J'}\Phi$, so taking the increasing union over $J'$ proves
\eqref{glue:crossing-quotient}. This calculation precedes any
closedness assertion. The crossing count is finite by
Proposition~\ref{op:operator}.

The image $a(F_J\Phi)$ is compact by continuity. The quotient in
\eqref{glue:crossing-quotient} is finite-dimensional over the finite
field $k$, hence is a finite set. Thus $a\Phi\cap F_J\Phi$ is a
finite union of translates of $a(F_J\Phi)$, and is compact and
closed. The $F_J\Phi$ form an open cover of $\Phi$, so $a\Phi$ is
closed in $\Phi$. The quotient is Hausdorff, and the quotient map
is open, proving the asserted compact open subspaces of $\mathcal Q$.
The kernel is closed by continuity, and its intersections with the
compact open stages give the assertion for $\mathcal K$.
\end{proof}

\subsection{The complete cone and its overlap}

In the category of complete objects, or equivalently in the perfect
$\mathscr E$-module calculus of
Proposition~\ref{complete:realization}, define
\begin{equation}
 m=\cofib(a:\widehat T'\longrightarrow\widehat T'),
 \qquad Y=\bigoplus_{r=0}^{D-1}m[r].
 \label{glue:complete-object}
\end{equation}
Recall that $W_i(Y)=\pi_i LY$ denotes the normalized overlap.
The finite sum over all residues ensures that both kinds of defect
occur in every degree.

\begin{lemma}[Compact stages of the cone]
\label{glue:overlap}
For every degree $i$, the normalized overlap has an identification
\begin{equation}
 W_i(Y)=\mathcal Q\oplus\mathcal K,
 \label{glue:overlap-splitting}
\end{equation}
compatible with the $t$-period identifications. The map
$\pi_iY\longrightarrow W_i(Y)$ has finite kernel and compact open
image. Each such image is commensurable with
\begin{equation}
 C_0^{\rm lat}=\im_{\mathcal Q}(F_0\Phi)
                     \oplus(\mathcal K\cap F_0\Phi).
 \label{glue:standard-compact}
\end{equation}
The topology of $\pi_iY$ is the intrinsic compact topology from its
finite reductions, and the induced topology on the overlap is the
locally compact topology in \eqref{glue:overlap-splitting}. Every
fixed compact stage of this overlap is contained in the image of
one sufficiently late $t$-translated group $\pi_{i+sD}Y$.
\end{lemma}

\begin{proof}
The overlap of $\widehat T'$ is supported in degrees divisible by
$D$, with normalized coefficient space $\Phi$ and the action of $a$
used above. Since $D>2$, the cofiber calculation in each summand has
separate cokernel and kernel residues. In degree $i$, the cokernel
component comes from $m[r]$ with $i\equiv r\bmod D$, by the cofiber
inclusion from $\widehat T'[r]$. The kernel component comes from
$m[r']$ with $i\equiv r'+1\bmod D$, by the connecting map into
$\widehat T'[r'+1]$. Summing gives
\eqref{glue:overlap-splitting} and its periodic compatibility.

For complete homotopy, Proposition~\ref{complete:sections} replaces
$\Phi$ by the compact coefficient stages
\[
 (x^{-s}R)^\Lambda=F_{\lfloor s/q\rfloor}\Phi.
\]
The two components may use different values of $s$ and hence of $J$;
both increase without bound under positive $t$-translation.
On a cokernel component, the kernel of the map from the stage
cokernel to $\mathcal Q$ is exactly
$(a\Phi\cap F_J\Phi)/a(F_J\Phi)$, finite by
Lemma~\ref{glue:pivot-topology}; its image is the compact open image
of $F_J\Phi$. On a kernel component, the stage gives simply
$\mathcal K\cap F_J\Phi$, injecting into $\mathcal K$. This proves
the finite-kernel and compact-open claims. Any two compact open
subspaces here are commensurable: their intersection is open in each
compact space, so the two quotient spaces are finite.

By Lemma~\ref{complete:reductions}, the kernel and cokernel
calculations at a fixed compact stage are inverse limits of the
corresponding finite-level calculations. These inverse systems of
finite groups satisfy the Mittag--Leffler condition: at each fixed
level, their decreasing sequences of images stabilize. Their
inverse limits are therefore exact, and the resulting compact
topology is the intrinsic topology from reductions. Taking the
$t$-colimit recovers the same locally compact union topology on
$\mathcal Q\oplus\mathcal K$.
Finally the coefficient stages increase without bound under positive
$t$-translation. Choosing a single sufficiently large stage for both
components proves the last assertion.
\end{proof}

\subsection{A lattice pairing the defects}

We choose the open part by constructing, in each degree, a discrete
complement to the compact open subspace
$C_0^{\rm lat}\subset\mathcal Q\oplus\mathcal K$.
Such a complement can be obtained by lifting an algebraic basis of
$(\mathcal Q\oplus\mathcal K)/C_0^{\rm lat}$ and taking its span.
Adding a vector of $\mathcal K\cap F_0\Phi$ to the kernel coordinate
of a lift does not change its quotient class. We use this freedom to
attach the distinguished compact kernel vectors to independent polar
cokernel classes.

For each forever unmatched positive row $j$, let $q_j\in\mathcal Q$
be the class of $w^j$. For each forever unmatched nonpositive column
$l$, choose $k_l\in\mathcal K$ with leading term $w^l$, as supplied
by Lemma~\ref{glue:pivot-topology}. In particular
$k_l\in\mathcal K\cap F_0\Phi\subset R$. Pair these lists in order
of increasing $j$ and increasing $|l|$; write $(j,l)$ for a resulting
pair of an unmatched row and an unmatched column.
Since $(0,k_l)\in C_0^{\rm lat}$, the pair $(q_j,k_l)$ and the
vector $(q_j,0)$ have the same image in the quotient. The next lemma
verifies the independence of these quotient classes and records
simultaneous bounds on their leading positions.

\begin{lemma}[Simultaneous defect pairs]
\label{glue:pairs}
The vectors $(q_j,k_l)\in\mathcal Q\oplus\mathcal K$, indexed by
the paired lists, are linearly independent modulo $C_0^{\rm lat}$.
For every sufficiently large integer $N$, at least $cN^\beta$ pairs,
for a fixed $c>0$, simultaneously satisfy $j\le N$ and $|l|\le N$.
For these pairs, their kernel coordinates are linearly independent
modulo $x^{qN+1}R$.
\end{lemma}

\begin{proof}
The $q_j$ are independent modulo $\im_{\mathcal Q}(F_0\Phi)$.
Otherwise a finite nontrivial relation, after subtracting a vector in
$F_0\Phi$, would be an actual image of $a$ whose leading position
is an unmatched positive row. This contradicts
Lemma~\ref{glue:pivot-topology}. It also proves independence of the
pairs modulo $C_0^{\rm lat}$.

Let $n_+(N)$ count the positive unmatched rows at most $N$, and let
$n_-(N)$ count the nonpositive unmatched columns of absolute value
at most $N$. In the ordered pairing, the first condition holds
precisely for the first $n_+(N)$ pairs and the second for the first
$n_-(N)$ pairs. Thus exactly $\min\{n_+(N),n_-(N)\}$ pairs satisfy both
conditions. Proposition~\ref{op:operator} bounds each count below
by a constant times $N^\beta$ for every sufficiently large $N$,
giving the simultaneous bound.
The leading terms of their kernel vectors are the distinct powers
$w^l=x^{q|l|}$, with $q|l|\le qN$. A nonzero combination has one of
these as its lowest $x$-power, so cannot lie in $x^{qN+1}R$.
\end{proof}

In every period residue, extend the independent images of the pairs
to an algebraic basis of
$(\mathcal Q\oplus\mathcal K)/C_0^{\rm lat}$. Choose lifts retaining
the specified pairs, and let $\mathcal L_i$ be their span. The
projection of this algebraic span to the quotient is an isomorphism.
It therefore meets $C_0^{\rm lat}$ trivially and, together with
$C_0^{\rm lat}$, fills the whole space. Since $C_0^{\rm lat}$ is
open, the span is discrete. Each open coset of $C_0^{\rm lat}$
contains exactly one point of the span, so the span is also closed.
Finally, the compact space $C_0^{\rm lat}$ maps onto the quotient
by the span, proving cocompactness. Make the same choice along each
$t$-period.

By Lemma~\ref{complete:open-splitting}, the $t$-open category is
the split category of $D$-periodic complexes of $k$-vector spaces.
There is consequently an open object $V$, also regarded as pushed
forward to $\mathbf D(\mathcal N)$, and a map $V\to LY$ whose maps
on homotopy are precisely these lattice inclusions. Define
\begin{equation}
 U=Y\mathop{\times}_{LY}V\quad\text{in }\mathbf D(\mathcal N),
 \label{glue:definition-U}
\end{equation}
where the pullback is a homotopy pullback.

\begin{proposition}[The glued object]
\label{glue:object}
The object $U$ of \eqref{glue:definition-U} satisfies
\[
 \pi_iU\text{ is finite for every }i,\qquad
 \widehat U\simeq Y\in\Thick(\widehat T'),\qquad LU\simeq V.
\]
Under these equivalences the structure map
$LU\longrightarrow L\widehat U$ is the chosen lattice map.
\end{proposition}

\begin{proof}
By Lemma~\ref{glue:overlap}, the map from $\pi_iY$ to $\pi_iLY$
has finite kernel and compact open image; call this image $C_i$.
It is commensurable with $C_0^{\rm lat}$, and
$\pi_iV=\mathcal L_i$ is a discrete complement to
$C_0^{\rm lat}$. Consider the joint map
\[
 \pi_iY\oplus\pi_iV\longrightarrow\pi_iLY.
\]
After dividing its kernel by the finite kernel of $\pi_iY\to C_i$,
one obtains $C_i\cap\mathcal L_i$. This intersection injects into
$C_i/(C_i\cap C_0^{\rm lat})$, which is finite. The cokernel of the
joint map is a quotient of
$C_0^{\rm lat}/(C_i\cap C_0^{\rm lat})$, also finite.
Thus the joint map has finite kernel and cokernel in every degree.
The pullback long exact sequence expresses $\pi_iU$ as an extension
of the kernel in degree $i$ by the cokernel in degree $i+1$.
Both are finite, so $\pi_iU$ is finite.

Both $V$ and $LY$ are $t$-invertible and hence have zero completion.
Exactness of completion and completeness of $Y$ give
$\widehat U\simeq Y$. Localizing the pullback instead gives
$LU\simeq V$, since the localization of $Y\to LY$ is an equivalence.
These identifications preserve the structure map. The object $Y$ is
a finite sum of shifts of a two-cell cone on $\widehat T'$, so lies
in $\Thick(\widehat T')$.
\end{proof}

The construction has now supplied the finite homotopy and complete
thick-membership parts of Proposition~\ref{transfer:criterion}.
The following consequence of the chosen lattice is the restriction
that will obstruct global thick membership.

\begin{lemma}[Constraint on global maps]
\label{glue:global-maps}
For every map $g:U\to P$ with $P\in\Thick(T')$, the induced map
\[
 W_i(\widehat g):\mathcal Q\oplus\mathcal K
                    \longrightarrow W_i(\widehat P)
\]
is continuous and sends each paired vector $(q_j,k_l)$ into the
discrete lattice $\im(\pi_iLP\to W_i(\widehat P))$.
\end{lemma}

\begin{proof}
Objects of $\Thick(T')$ have finite homotopy in each degree. The
same is true of $U$ by Proposition~\ref{glue:object}.
Proposition~\ref{complete:topology} therefore makes every map between
their completions continuous on overlap homotopy. Under
Lemma~\ref{glue:overlap}, the source topology is precisely the
locally compact topology on $\mathcal Q\oplus\mathcal K$ used here.

The additional lattice assertion uses that $\widehat g$ comes from
the global map $g$. Naturality of the fracture square gives a
commutative square from $LU\to L\widehat U$ to
$LP\to L\widehat P$. Each paired vector is represented in
$\pi_iLU=\pi_iV=\mathcal L_i$. Applying $\pi_iLg$ and then the
target structure map places its image in
$\im(\pi_iLP\to W_i(\widehat P))$, which is the discrete lattice
provided by Proposition~\ref{complete:topology}. This argument
uses the representative in $\pi_iLU$; an individual pair need not
lift to $\pi_iU$.
\end{proof}

To see how the pairing will be used, choose a compact open target
subspace containing the images of the kernel coordinates. Modulo
this subspace, each paired image is its polar image. Since all paired
images lie in a discrete lattice, the part of their span inside the
compact subspace has finite dimension. Thus a rank bound for the
polar images controls the rank of the paired images up to a fixed
finite error. The next section obtains the required bound under its
homology-vanishing hypothesis.
